% §VIII --- Limitations (~0.25 page, ~250 words; >=6 honest items).
% Drafted D11 per ml-paper-writing checklist (Limitations are required,
% reviewers are explicitly told not to penalize honesty).
\section{Limitations}\label{sec:limit}

We list six concrete limitations that we did not fully resolve. First,
the system model assumes a single base station co-located with the MEC
server. Multi-cell extensions with handover and inter-cell interference
are out of scope and would require both a new association rule and a
hierarchical critic. Second, every UAV in our experiments flies at a
fixed altitude $H_U=50$~m with the simplified flight / hover energy
constants of \cite{liu2024jamminguav}; full 3-D trajectory control with
realistic aerodynamics is a natural extension that requires a richer
continuous action space. Third, the DeepSC-VQA semantic similarity
look-up table that drives \eqref{eq:simlut} is discretized on a
$5\!\times\!5$ symbol grid and a $7$-bin SINR axis; whether \method's
gains transfer to LUTs derived from different semantic encoders or
different VQA datasets remains unverified, and we leave this as a
priority follow-up. Fourth, both \method and the MAPPO-hybrid baseline
share the same training sample budget; whether the role-aware actor
remains preferable in the deeply sample-constrained regime
($\le 10^{4}$ rollouts) is open. Fifth, \method requires substantially
more training compute than classical convex-optimization-based
allocation methods; the wall-clock numbers in
Table~\ref{tab:wallclock} make the trade-off explicit but reviewers
deploying in compute-constrained settings should weigh this carefully.
Finally, the adversarial-jamming experiments in
Fig.~\ref{fig:jam} use a constant-power white-noise attacker; learning
adversaries that adapt to \method's policy in a min--max framework are
not modeled and may degrade the reported robustness margins.
